%!TEX root = ../dissertation.tex

Composed image retrieval (CIR) searches an image database with a query made of a
reference image and a short modification text that states how the desired result
should differ from it. The recent i-CIR benchmark makes the task markedly harder
by defining relevance at the level of a specific object \emph{instance} and by
surrounding every query with hard negatives that defeat any method relying on a
single modality. Its accompanying baseline, \textsc{Basic}, shows that frozen
vision--language features combined with careful, training-free post-processing
are surprisingly competitive; yet it treats the modification text in isolation,
even though a phrase such as ``make it red'' is meaningless to a text encoder
without knowing \emph{what} should become red.

This thesis uses language as an explicit bridge between the two modalities.
A frozen multimodal large language model (InternVL-3.5-8B) reads the reference
image and the modification \emph{together} and writes a short description of the
\emph{imagined target} -- effectively performing the edit in language space. This
self-contained \emph{target caption} is embedded by the frozen text encoder and
fused with the image and text similarities through a clamped triplet product that
acts as a soft logical \textsc{and}: a candidate must agree with all three signals
to rank highly. The pipeline is entirely zero-shot -- every encoder, the caption
generator, and all fusion and post-processing steps are off-the-shelf and frozen.
Caption variance is reduced by averaging the embeddings of five complementary
instruction prompts (\texttt{avg-5}).

The method is developed and ablated on a $20\%$ instance subset of i-CIR and the
successful configurations are then re-run at full scale ($1{,}883$ queries against
$752{,}092$ database images) with four backbones (CLIP ViT-L/14, CLIP ViT-H/14,
SigLIP ViT-L/16, and SigLIP2 g/16-384). The target caption is the single most
effective component: it improves every backbone, and on the best backbone it
contributes a larger gain than the entire \textsc{Basic} post-processing stack it
is designed to replace. A plain text$\times$image baseline augmented with the
caption reaches, and with a strong backbone exceeds, the fully tuned
\textsc{Basic} pipeline. The caption and feature centering are complementary, and
the best backbone within the full pipeline is SigLIP2 -- not the strongest
bare-baseline encoder -- reaching $61.98$ macro-mAP, the best result of the thesis.
A token-level failure analysis shows that the remaining errors are concentrated in
modifications whose targets are absent from the database (state changes, domain
transfers, out-of-context wear) rather than in the retrieval mechanism itself.
